\section{Methodology}
\label{meth}
The spatial and temporal distribution of global air traffic is highly inhomogeneous, due to variations in passenger demand, safety regulations, airspace organisation and flight route optimisation \cite{Gossling2020}. Due to the relative unpredictability and non-uniformity in air traffic flows and patterns, modelling and investigation should ideally be based around real-world data. In this paper, historical commercial aircraft trajectory data obtained from the open-source air traffic dataset OpenSky \cite{Schafer2014}, is leveraged to explore how aircraft plumes typically overlap in airspace regions where air traffic densities are particularly high. The OpenSky Network sources data from Automatic Dependent Surveillance-Broadcast (ADS-B) ground-based receivers, which determine aircraft latitude, longitude and altitude, along with additional measurements such as flight track and ground speed, at a temporal resolution of one second. This dataset is used here to deduce information regarding air traffic flows and airspace density, as well as modelling of aircraft plume intersections, throughout the month of January 2022. All necessary computations were carried out using the Python programming language, coupled with the traffic API \cite{Olive2019} to handle OpenSky data requests and all plots were generated using Python packages Matplotlib and Cartopy.

In flight air traffic sectors where demand is considerably high, or where two or more flight corridors intersect, air traffic can become significantly congested, with airspace capacities often approaching their upper limit \cite{Odoni1987}. Although undesirable from an air traffic safety and efficiency perspective, it so happens that these regions could provide excellent insight into the real-world practicality of formation flight and controlled plume overlap. Therefore, the first aim of study was to use OpenSky ADS-B data to locate the highest density flight regions during a typical busy time period. January 2022 was selected, as this represented a point in time when the aviation sector had recovered significantly from the COVID downturn, and was prior to the Russian invasion of Ukraine, which further disrupted the air traffic situation considerably. Trajectory data over this month were discretised into 1\textdegree\ longitude $\times$ 1\textdegree\ latitude bins, at a typical cruise altitude of 38,000 ft $\pm$500~ft. The categorised data were then used to construct a global heatmap of mean airspace density for each grid cell. To reduce computational intensity, snapshots of the air traffic situation were taken at the start of every hour, each day of the month. 

Following the outcome of the global analysis, two key high-density hotspots were selected for further analysis. The flux densities through the box were calculated for each day and averaged over the whole month, to single out the time steps with the highest average aircraft count. The case study of these two hotspots continued at the subgrid scale, through higher resolution observation of trajectories and modelling of their associated exhaust plumes. Dispersion modelling of aircraft plumes has been covered extensively in the literature over the past few decades, with models of varying fidelities becoming widely accepted amongst the community \cite{Paoli2011}. For this paper, the aim was to select a relatively low fidelity model to track the plume dispersion from OpenSky aircraft trajectories, whilst also being computationally efficient, so that many trajectories could be easily tested. The model used here was that of the Multi-Shell Plume model from Kraab{\o}l et al. \cite{Kraabol2000a}. Here, the plume's horizontal, vertical and skewed standard deviations ($\sigma_{\mathrm{h}}$, $\sigma_{\mathrm{v}}$ and $\sigma_{\mathrm{s}}$) and radii ($r_{\mathrm{h}}$, $r_{\mathrm{v}}$) of the outermost shell were modelled through equations (3.1) to (3.5). 

\begin{align}
\begin{split}
	\label{sigma_h}
	{\sigma_{\mathrm{h}}}^{2}(t) {}&= \frac{2}{3} s^2 D_{\mathrm{v}} t^3 + (2 D_{\mathrm{s}} + s {\sigma_{\mathrm{0v}}}^2) s t^2 + 2 D_{\mathrm{h}} t + {\sigma_{\mathrm{0h}}}^2,
\end{split} \\
\begin{split}
	{\sigma_{\mathrm{v}}}^{2}(t) {}&= 2 D_{\mathrm{v}} t + {\sigma_{\mathrm{0v}}}^2,
	\end{split}\\
	{\sigma_{\mathrm{s}}}^{2}(t) {}&= s D_{\mathrm{v}} t^2 + (2 D_{\mathrm{s}} + s {\sigma_{\mathrm{0v}}}^2) t. \\[10pt]
	r_h {}&= 3 \sigma_{\mathrm{h}} \\
	\label{r_v}
	r_v {}&= 3 \sigma_{\mathrm{v}}
\end{align}

These equations derive from the two-dimensional diffusion equation, and hence contain initialised atmospheric parameters, according to the assumed turbulent conditions of Durbeck and Gerz \cite{Durbeck1996}: wind shear (s $=$ 0.005~\ce{s^{-1}}), horizontal (\ce{D_h} $=$ 16.7~\ce{m^2 s^{-1}}), vertical (\ce{D_v} $=$ 1.1 \ce{m^2 s^{-1}}) and shear (\ce{D_s} $=$ 0.5(\ce{D_v} \ce{D_h})\ce{^2}) diffusion coefficients, at an arbitrary time step t. The time step used in the plume simulations present throughout this paper is 10~s. This means that aircraft positional data were gathered every 10~s, and instances of plume segments were initialised, for active aircraft in the grid cell at that particular time step. For plume instances that already exist, they are expanded each time step, according to the above equations.

In addition to plume modelling, data were collected on aircraft-to-aircraft information across all flights that traverse through the observed grid cells over the specified timeframe. This includes relative distance, relative track and bearing between aircraft pairs. From this, criteria for scenarios where it was likely that two plume instances would overlap was determined, thus promoting climate-beneficial saturation effects, as elaborated on in section \ref{Intro}. The criteria to be met to achieve successful plume overlap were termed \textit{significant plume encounters} (SPEs). Based on the E-folding time of \ce{NO_x} and typical contrail growth rates at aircraft cruising altitudes, it was decided that a threshold overlap time for an SPE was around 2 hours \cite{Shallcross2022, Kraabol2000b}. Thus, the final objective of the study was to determine how many instances of aircraft pairs, at each time step, achieved a successful SPE throughout the observed duration. 
\vspace{-.25cm}
\begin{figure}[H]
\centering
		\includegraphics[width=0.70\textwidth]{3/Plume_overlap.pdf}
		%\vspace{-.1cm}
		\caption{Schematic drawing of SPE criteria used to deduce plume overlap cases and compare the relative importance of plume intersection events.}
		\label{SPE}
\end{figure}

SPEs were ranked by significance from 1 to 4, as is visually represented in figure \ref{SPE}. Under assumed plume conditions, horizontal plume radius \ce{r_h} is predicted to reach between 6 and 7~km at the 2~h threshold, whilst \ce{r_v} approaches a mere 75--100~m. SPE ranks were designed in such a way as to capture all instances where plumes from aircraft pairs are likely to overlap, prior to the proposed 2~h threshold. \ce{SPE_1} is classed as any aircraft that is within a 20~km proximity of the baseline aircraft, has a bearing within $\pm$ 45\textdegree \ of the baseline track, and also, has a track that is within $\pm$ 45\textdegree \ of the baseline track. For \ce{SPE_2}, the same rules on relative tracks apply, however relative distance is between 20 and 50~km. \ce{SPE_3} counts as any aircraft that is outside the $\pm$ 45\textdegree\ boundary, however is within 20~km. Relative headings were neglected here, as any aircraft crossing at such obtuse angles will likely experience plume overlap to some degree. Finally, \ce{SPE_4} represents any aircraft that does not satisfy any of the other conditions, however is still within a 50~km radius of the baseline aircraft. The ranks were deduced based on relative distance and tracks between all aircraft pairs at every time step, prioritising tracks over distance. This is because two aircraft flying a similar heading are extremely likely to experience significant and persistent plume overlap, even if they are a large distance apart longitudinally. This is because aircraft travel at high speeds and trajectories tend not to deviate too far from their heading, during the cruise segment of flight \cite{Tait2022}. 

% Which altitude column you’ve selected the data using (GPS or Barometric – both have pros and cons).


% You have all the detail on your grids etc. – I’ve tried digging to find specific info about air corridors in that region, its all rather complicated! Though its clear high altitude airways skirt around the edge of France, leaving the centre clearish of high altitude airways